\documentclass[11pt, a4paper]{article}
\usepackage[utf8]{inputenc}
\usepackage[swedish]{babel}
\usepackage{amsmath, amssymb, amsthm}
\usepackage{geometry}
\usepackage{hyperref}

\geometry{top=2.5cm, bottom=2.5cm, left=2.5cm, right=2.5cm}

\title{\textbf{Hvitfeldtska Spetsutbildning Matematik: Kunskapsöversikt}}
\author{Förberedd för math.bluehawana.com}
\date{\today}

\begin{document}

\maketitle

\section{Introduktion}
Detta dokument sammanfattar de centrala matematiska områden och begrepp som krävs för att lyckas på antagningsprovet till Hvitfeldtskas spetsutbildning i matematik.

\section{Aritmetik och Talteori}
\begin{itemize}
    \item \textbf{Räkneordning:} Prioriteringsregler (parenteser, potenser, multiplikation/division, addition/subtraktion).
    \item \textbf{Primtal och Delbarhet:}
    \begin{itemize}
        \item Primtalsfaktorisering.
        \item Delbarhetsregler (2, 3, 4, 5, 6, 8, 9).
        \item Största gemensamma delare (SGD) och minsta gemensamma multipel (MGM).
    \end{itemize}
    \item \textbf{Bråk och Procent:}
    \begin{itemize}
        \item Förlängning och förkortning av bråk.
        \item Omvandling mellan bråk, decimalform och procent.
        \item Förändringsfaktor.
    \end{itemize}
    \item \textbf{Potenser och Rötter:}
    \begin{itemize}
        \item Potenslagarna ($a^x \cdot a^y = a^{x+y}$, $(a^x)^y = a^{xy}$, etc.).
        \item Kvadratrötter och högre rötter.
        \item Grundläggande uppskattning av rötter (t.ex. $\sqrt{50} \approx 7.07$).
    \end{itemize}
\end{itemize}

\section{Algebra och Funktioner}
\begin{itemize}
    \item \textbf{Algebraiska uttryck:} Förenkling, parenteshantering, faktorisering.
    \item \textbf{Ekvationer:}
    \begin{itemize}
        \item Linjära ekvationer.
        \item Ekvationssystem (substitutionsmetoden, additionsmetoden).
        \item Andragradsekvationer (pq-formeln, nollproduktmetoden).
    \end{itemize}
    \item \textbf{Olikheter:} Hantering av olikheter, intervall.
    \item \textbf{Mönster och Talföljder:}
    \begin{itemize}
        \item Aritmetiska och geometriska talföljder.
        \item Rekursiva formler.
    \end{itemize}
\end{itemize}

\section{Geometri}
\begin{itemize}
    \item \textbf{Plangeometri:}
    \begin{itemize}
        \item Vinklar och vinkelsummor (triangel $180^\circ$, fyrhörning $360^\circ$).
        \item Area och omkrets av trianglar, rektanglar, parallelltrapetser, cirklar.
        \item Pythagoras sats.
    \end{itemize}
    \item \textbf{Likformighet och Kongruens:}
    \begin{itemize}
        \item Topptriangelsatsen och transversalsatsen.
        \item Skala (längdskala, areaskala, volymskala).
    \end{itemize}
    \item \textbf{Cirkelgeometri:}
    \begin{itemize}
        \item Randvinkelsatsen.
        \item Kordasatsen.
        \item Tangenter till cirklar.
    \end{itemize}
\end{itemize}

\section{Kombinatorik och Sannolikhet}
\begin{itemize}
    \item \textbf{Kombinatorik:}
    \begin{itemize}
        \item Multiplikationsprincipen.
        \item Permutationer (ordning spelar roll).
        \item Kombinationer (urval utan ordning).
    \end{itemize}
    \item \textbf{Sannolikhet:}
    \begin{itemize}
        \item Klassisk sannolikhetsdefinition ($P(A) = \frac{\text{gynnsamma utfall}}{\text{möjliga utfall}}$).
        \item Oberoende händelser.
        \item Komplementhändelse.
    \end{itemize}
\end{itemize}

\section{Problemlösningsstrategier}
\begin{itemize}
    \item \textbf{Rita en figur:} Alltid första steget i geometriproblem.
    \item \textbf{Testa små fall:} Förstå mönster genom att testa $n=1, 2, 3$.
    \item \textbf{Arbeta baklänges:} Ibland lättare att utgå från vad som ska visas.
    \item \textbf{Variabelinförande:} Beteckna sökta storheter med $x, y$.
\end{itemize}

\end{document}
